\section{Anhang: Nutzung der Simulation}

Die \href{https://github.com/danielschwalbe/ipres}{Simulation} ist ein in der
\href{https://www.python.org}{Python}-Programmiersprache geschriebenes Projekt. Es enthält neben
der \textbf{ipres} Python-Package viele Readmes und Beispielnotebooks.
\href{https://jupyter.org}{Jupyter Notebooks} ist eine aus der Python-Welt stammende Technologie,
die es erlaubt, Markdown (Dokumentation) und Code in einer Datei zusammenzufassen und interaktiv
auszuführen. Das macht sie geeignet, um z.\,B.\ eben solche Wahlverfahren zu demonstrieren. Man
kann Parameter ändern und live sehen, wie sie sich auswirken.

Eigentlich müsste man einen Jupyter-Lab-Server starten, um die Notebooks auszuführen. Das Projekt
verlinkt aber zu \href{https://mybinder.org/v2/gh/danielschwalbe/ipres/HEAD}{binder}, welches das
Projekt in eine virtuelle Maschine kopiert und dort den Jupyter-Server bereits ausführt. Man muss
nur ein bisschen Geduld haben, bis die Maschine gestartet ist. Bis dahin sieht es so aus:

\begin{figure}[h]
  \centering
  % TODO: Screenshot einfügen
  % \includegraphics[width=0.85\textwidth]{screenshots/binder_start.png}
  \fbox{\parbox{0.85\textwidth}{\centering\vspace{2cm}
    \textit{[Screenshot: Binder-Startseite beim Laden der virtuellen Maschine]}
    \vspace{2cm}}}
  \caption{Binder lädt die virtuelle Maschine}
\end{figure}

Wenn die Maschine hochgefahren ist, wird man von folgendem Bildschirm begrüßt:

\begin{figure}[h]
  \centering
  % TODO: Screenshot einfügen
  % \includegraphics[width=0.85\textwidth]{screenshots/jupyterlab_start.png}
  \fbox{\parbox{0.85\textwidth}{\centering\vspace{2cm}
    \textit{[Screenshot: JupyterLab-Startseite nach dem Hochfahren]}
    \vspace{2cm}}}
  \caption{JupyterLab nach dem Start}
\end{figure}

Um zu den Notebooks zu gelangen, muss man auf \texttt{notebooks} klicken.

\begin{figure}[h]
  \centering
  % TODO: Screenshot einfügen
  % \includegraphics[width=0.85\textwidth]{screenshots/notebooks_folder.png}
  \fbox{\parbox{0.85\textwidth}{\centering\vspace{2cm}
    \textit{[Screenshot: Notebooks-Ordner mit \texttt{de/} und \texttt{en/} Unterordnern]}
    \vspace{2cm}}}
  \caption{Notebooks-Ordner}
\end{figure}

Und dann auf \texttt{de}, um zu den deutschen Notebooks zu gelangen.

\begin{figure}[h]
  \centering
  % TODO: Screenshot einfügen
  % \includegraphics[width=0.85\textwidth]{screenshots/notebooks_de.png}
  \fbox{\parbox{0.85\textwidth}{\centering\vspace{2cm}
    \textit{[Screenshot: Inhalt des \texttt{notebooks/de/}-Ordners]}
    \vspace{2cm}}}
  \caption{Deutsche Notebooks}
\end{figure}

Hier kann man sich dann welche herauspicken, z.\,B.\ die \texttt{bundestagswahl.ipynb}.

Durch Klick auf den Startpfeil wird die gerade markierte Zelle ausgeführt. Man beachte, dass die
Zellen in der Regel von oben nach unten ausgeführt werden müssen, weil die unteren Zellen
meistens (nicht immer) die Ergebnisse der oberen Zellen voraussetzen.

\begin{figure}[h]
  \centering
  % TODO: Screenshot einfügen
  % \includegraphics[width=0.85\textwidth]{screenshots/notebook_execution.png}
  \fbox{\parbox{0.85\textwidth}{\centering\vspace{2cm}
    \textit{[Screenshot: Ausführung einer Notebook-Zelle mit Startpfeil]}
    \vspace{2cm}}}
  \caption{Ausführung einer Notebook-Zelle}
\end{figure}

Mithilfe der zwei Pfeile kann man auch das ganze Notebook auf einmal ausführen.

\begin{figure}[h]
  \centering
  % TODO: Screenshot einfügen
  % \includegraphics[width=0.85\textwidth]{screenshots/notebook_run_all.png}
  \fbox{\parbox{0.85\textwidth}{\centering\vspace{2cm}
    \textit{[Screenshot: \enquote{Run All}-Schaltfläche zum Ausführen aller Zellen]}
    \vspace{2cm}}}
  \caption{Alle Zellen auf einmal ausführen}
\end{figure}

Wenn man am Code etwas ändern will, kann man die Zeilen einfach editieren und nochmal ausführen.
Man könnte z.\,B.\ in der oberen Zelle andere Seeds eintragen und würde sehen, wie andere Parteien
gewinnen.

Falls man weitergehende Änderungen machen möchte, muss man deswegen heutzutage nicht unbedingt
Python können. Die KIs können das auch und helfen gerne.
